\documentclass[11pt]{article}
% Shared preamble for per-Python-file documentation (input via \input{...}).
\usepackage[margin=1in]{geometry}
\usepackage[T1]{fontenc}
\usepackage[utf8]{inputenc}
\usepackage{lmodern}
\usepackage{hyperref}
\usepackage{xcolor}
\usepackage{listings}
\usepackage{listingsutf8}
\lstset{
  basicstyle=\ttfamily\small,
  columns=fullflexible,
  breaklines=true,
  upquote=true,
  inputencoding=utf8,
  extendedchars=true,
  frame=single,
  framerule=0.2pt,
  tabsize=2
}


\title{\texttt{model/d2q.py} (Q\&A)}
\author{}
\date{}

\begin{document}
\maketitle

\paragraph{Q: What is the purpose of \texttt{model/d2q.py}?}
\textbf{A:} It implements the D2Q baseline model, which predicts a scalar ``quantile'' score (after sigmoid) from features. The accompanying run script (\lstinline|run_d2q.py|) maps watch-time labels into bucket-specific quantile space during training and maps predictions back to watch-time values during evaluation.

\paragraph{Q: What are the inputs and outputs of \lstinline|D2Q.forward|?}
\textbf{A:} Inputs are the standard feature dictionary \lstinline|x_dict|; outputs are a scalar per example (a sigmoid-activated value) intended to represent a normalized quantile within a duration bucket.

\paragraph{Q: How are features encoded?}
\textbf{A:} It uses embeddings for sparse and sequence features and linear transforms for continuous features, concatenates them, and feeds the result into \lstinline|MultiLayerPerceptronD2Q| (a Swish-based MLP).

\paragraph{Q: Why does D2Q need duration buckets?}
\textbf{A:} D2Q is designed to model the conditional distribution of play time given video duration by first bucketing duration. The mapping from watch time to quantile is bucket-specific (precomputed by preprocessing), and the model predicts in that normalized space.

\paragraph{Q: Code example: forward pass shape (schematic)?}
\textbf{A:}
\begin{lstlisting}[language=Python]
q = model(features)   # q: [batch] after sigmoid
\end{lstlisting}

\paragraph{Q: Full source code?}
\textbf{A:}
\lstinputlisting[language=Python]{/workspace/model/d2q.py}

\end{document}
